\documentclass [xcolor=svgnames, t] {beamer} 
\usepackage[utf8]{inputenc}
\usepackage{booktabs, comment} 
\usepackage[absolute, overlay]{textpos} 
\usepackage{pgfpages}
\usepackage[font=footnotesize]{caption}
\useoutertheme{infolines} 

\definecolor{gold}{RGB}{254, 206, 0}

%\setbeamercolor{title in head/foot}{bg=gold, fg=black}
\setbeamercolor{author in head/foot}{bg=myuniversity}
\setbeamertemplate{page number in head/foot}{}
\usepackage{csquotes}
\setbeamertemplate{navigation symbols}{}

\usepackage{amsmath}
\usepackage[makeroom]{cancel}


\usepackage{textpos}

\usepackage{tikz}
\usepackage{multirow}

\usetheme{Dresden}
\setbeamertemplate{footline}[]
\definecolor{myuniversity}{RGB}{0, 85, 150}
\usecolortheme[named=myuniversity]{structure}
\usepackage{tikz}



\title[]{A New Formulation of Minimum Risk Fixed-Width Confidence
Interval (MRFWCI) for a Normal Mean}
\titlegraphic{\includegraphics[height=1 cm]{logo.png}}
\institute[]{Department of Statistics \\ University of Connecticut}

\author[]{ Swathi Venkatesan }

\date{}



\addtobeamertemplate{navigation symbols}{}{%
    \usebeamerfont{footline}%
    \usebeamercolor[fg]{footline}%
    \hspace{1em}%
    \insertframenumber/\inserttotalframenumber
}

\begin{document}
\begin{frame}
\maketitle
\end{frame}




%%%%%%%%%%%%%%%%%%%%%%%%%%

%\begin{frame}{Abstract}
%\vspace{10mm}
%The fixed-width confidence interval (FWCI) estimation problems for a normal mean when 
%the variance is unknown have moved along under a zero-one loss without taking into account sampling cost. 
%While, minimum risk point estimation (MRPE) problems have grown largely under squared error loss (SEL) plus sampling cost. 
%Here, a new formulation combining both MRPE and FWCI methodologies is introduced with desired asymptotic second-order 
%characteristics under a unified structure to develop a minimum risk fixed-width confidence interval (MRFWCI) strategy.
%\end{frame}

%%%%%%%%%%%%%%%%%%%%%%%%%%
\section{Motivation}
\begin{frame}{Motivation}

Under a classical fixed-width confidence interval (FWCI) methodology, one treats the associated confidence coefficient as the main focus, 
which is naturally supposed to go up as the sample size increases. 
In such a formulation, however, one does not encounter the cost of unit observation at all,
which often creates a false pretense that perhaps (i) observations may not cost at all
and/or (ii) one’s available budget may be unlimited.\\
\vspace{0.2cm}
On the other hand, under the minimum risk point estimation (MRPE) strategy, one openly balances the estimation
error by taking into account the cost of sampling, which go together hand in hand.\\
\vspace{0.2cm}
Here, we propose to formulate a new way to combine both MRPE and FWCI methodologies
with desired asymptotic second-order characteristics under one unified theory.

\end{frame}

%%%%%%%%%%%%%%%%%%%%%%%%%%


\section{Loss and Risk Function}
\begin{frame}{Construction of Loss and Risk Functions}
    Suppose $X_1,X_2, \dots ,X_n \dots$ are \textit{independent 
    and identically distributed (i.i.d)} observations from a common 
    $N(\mu,\sigma^2)$ distribution where $\mu \in \mathbb{R}$, 
    $\sigma \in \mathbb{R}^{+}$, $\theta=(\mu,\sigma)$ and both $\mu,\sigma$ are unknown. \\
    \vspace{2mm}
    Having recorded a fixed number $n (\geq 2)$ of observations $X_1,X_2,\dots,X_n$, we denote:
    $$\text{Sample mean: }\bar{X}_n=n^{-1}\sum_{i=1}^nX_i$$
    $$\text{Sample variance: }S_n^2=(n-1)^{-1}\sum_
    {i=1}^n {(X_i-\bar{X}_n)^2}$$
    \vspace{0.1cm}
    We also denote \\
    $\phi(x)=(2\pi)^{-1/2}exp(-x^2/2)$ and $\Phi(x) = \int_{-\infty}^{x} \phi(y) dy $
    
    \end{frame}
%%%%%%%%%%%%%%%%%%%%%%%%%%

    
 \begin{frame}{}
    \vspace{7mm}

    
    Having pre-assigned numbers, $d(>0)$ and $0<\alpha<1$, the FWCI for $\mu$ is:
        $$\textbf{FWCI: } J_n \equiv \left[ \bar{X}_n \pm d\right],$$
        which  has the fixed width, $2d$. The confidence coefficient associated with $J_n$ is given by
        \begin{equation} \label{opt_n} 
        \begin{split}
        & \text{ with confidence coefficient }P_{\theta}\{\mu \in J_n\} = 2\Phi(n^{1/2}d/\sigma)-1  \ge 1-\alpha \\
        & \text{ when } n \text{ is the smallest integer } \ge z^2_{\alpha/2}\sigma^2 d^{-2} \equiv n^*.
        \end{split}
        \end{equation}
        $z_{\alpha/2}$ stands for the upper $50 \alpha \%$ point of a standard normal r
        andom variable; that is, $\Phi(z_{\alpha/2})=1-\frac{\alpha}{2}$

\end{frame}
%%%%%%%%%%%%%%%%%%%%%%%%%%

 \begin{frame}{}
  \vspace{7mm}

The expression of $n^*$ from (1) is referred to as the \textit{optimal fixed sample size} required to 
construct the FWCI, $J_n$, for $\mu$ had $\sigma$ been known. \\
\vspace{0.2cm}
We emphasize that the magnitude of $n^*$ remains unknown because $\sigma^2$ is 
assumed unknown.Dantzig (1940) showed that there 
did not exist any fixed-sample-size methodology to solve this problem.\\
\vspace{0.2cm}
Stein (1945,1949), followed by Anscombe (1952, 1953), Ray (1957), Chow and Robbins (1965), and
others, developed a number of pioneering multistage and purely sequential sampling
methodologies in order to address the FWCI problem.

\end{frame}
%%%%%%%%%%%%%%%%%%%%%%%%%%
\begin{frame}{}
    \vspace{7mm}
    Generically, we may initiate the following idea of a loss function in practice:
        \begin{equation} \label{loss_fn}
            L_n(\theta)  = EstErr_n(\theta)+c_n(\theta),
        \end{equation}
        where $EstErr_n(\theta)$ is the estimation error involving 
        $I \left[ {\tau(\theta) \not\in \text{FWCI};n} \right]$ plus the $c_n(\theta) \equiv \text{sampling cost}_n$\\
        \vspace{2mm}
    We combine (1) and (2) to express:
        \begin{equation}
        p_n(\theta ;J_n) \equiv E_{\theta} \{ {I[ \mu \not\in J_n;n]}\} =2\{ 1-\Phi(n^{1/2}d/\sigma) \}.
        \end{equation}
        which converges to $0$ as $n \rightarrow \infty$.\\
        \vspace{0.1cm}
                
\end{frame}
%%%%%%%%%%%%%%%%%%%%%%%%%%
\begin{frame}{}
\vspace{10mm}
Having fixed pre-assigned $d,\alpha,\rho$ and $\sigma$, we incorporate the 
following expression of the cost function:
        \begin{equation} \label{cost_fn}
        c_n(\theta) = 2(d^\rho \sigma)^{-1} \phi(z_{\alpha/2})n^{1/2}.
        \end{equation}
Then, in the spirits of (1) and (3)-(4), we construct the loss function (for fixed $0<\rho \le 1$) as:
        \begin{equation}
            \textbf{Loss: }L_n(\mu,J_n) \equiv d^{-\rho-1}I[\mu \not\in J_n;n]+c_n(\theta).
        \end{equation}

\end{frame}

%%%%%%%%%%%%%%%%%%%%%%%%%%
\begin{frame}{}
\vspace{10mm}

The corresponding Risk is given by:
    \begin{equation}
    \begin{split}
        \textbf{Risk: } R_n(\mu,J_n)  
        & \equiv E_{\theta}[L_n(\mu,J_n)] \\ 
        &= 2d^{-\rho-1}[1-\Phi(n^{1/2}d/\sigma)]+2(d^{\rho}\sigma)^{-1}\phi(z_{\alpha/2})n^{1/2} 
    \end{split}
    \end{equation}
We may pretend to treat $n$ as a continuous variable 
and differentiate the expression from (6), 
which is convex in $n$, with respect to $n$ and equate it to zero to claim:\\
    $-d^{-\rho-1}(d/\sigma)n^{-1/2}\phi(n^{1/2}d/\sigma)+
    2(d^\rho \sigma)^{-1}\phi(z_{\alpha/2})\frac{1}{2} n^{1/2}=0$\\
    \vspace{0.1cm}
    $\Rightarrow \phi(n^{1/2}d/\sigma)=\phi(z_{\alpha/2}) \Rightarrow  n \equiv n^*=z^2_{\alpha/2}\sigma^2d^{-2}$
    
\end{frame}
%%%%%%%%%%%%%%%%%%%%%%%%%%
\section{General Theory-MRFWCI}
\begin{frame}{General setup of MRFWCI-Normal mean Estimation}
\vspace{5mm}
Let $m(\geq 2)$ be an initial or pilot sample size, not involving $d$. In the case of 
each estimation strategy that will be discussed later ,  we will readily see that the boundary 
conditions for stopping will depend only upon the sample variance(s), that is the event $[N=n]$ 
would only depend on $(S_m^2,\dots,S_n^2)$ for all $n \geq m$, ultimately arriving at the terminally accrued data, 
$\{N,X_1,\dots,X_N\}$ where $N$ would estimate $n^*$. \\
\vspace{0.4cm}
By using Basu's theorem, $I(N=n)$ will be independent of $\bar{X}_n$ for all $n \geq m$. \\
\vspace{0.1cm}
\end{frame}
%%%%%%%%%%%%%%%%%%%%%%%%%%

\begin{frame}
\vspace{10mm}
Now, we propose the MRFWCI for $\mu$ based on the recorded observation $\{N,X_1,X_2,\dots ,X_N\}$ as:
    \begin{equation*}
    \textbf{MRFWCI: }J_N=[\bar{X}_N \pm d] \text{ where }\bar{X}_N=N^{-1} \sum_{i=1}^N{X_i}.    
    \end{equation*}
    
    The risk function $R_N(\mu,J_N)$ with the terminal strategy $(N,J_N)$ for $\mu$ under (5) is:
    \begin{equation*}
        R_N(\mu,J_N) = 2\{ d^{-\rho-1} E_{\theta}[1-\Phi(N^{1/2}d/\sigma)]+ (d^{\rho}\sigma)^{-1}\phi(z_{\alpha/2})E_{\theta}(N^{1/2})\}.
    \end{equation*}
    
    %The optimal fixed sample size (had $\sigma$ been known) corresponding to the minimum risk associated with the MRFWCI 
    %formulation under the new loss function remains exactly the same as that shown earlier in (1).
    
\end{frame}

%%%%%%%%%%%%%%%%%%%%%%%%%%

%%%%%%%%%%%%%%%%%%%%%%%%%%
%\section{Properties}
\begin{frame}{}
\vspace{15mm}
 Without specifying a formal sampling strategy , we assume a number of key associated properties as $d \xrightarrow{} 0$ :
   
    \begin{itemize}
        \item[\textbf{A1}] $ E_{\theta}(N)=n^* +a_1 +o(1)$ [\textit{asymptotic second-order efficiency}];
         \item[\textbf{A2}] $H \equiv \frac{N-n^*}{n^{*1/2}} \xrightarrow{\L} N(0,a_2), a_2>0$ ;
          \item[\textbf{A3}] $P_{\theta}\{ N \le \epsilon^*\}=O(n^{*-a_3}),\epsilon \in (0,1), a_3>0$;
           \item[\textbf{A4}] $|H|^{a_4} \text{ is uniformly integrable },a_4 >0$.
    \end{itemize}
    
    where the expressions of the real numbers $a_i's$ would not involve $d$.\\
    \vspace{0.2cm}
    The asymptotic characteristics assumed above will hold under appropriate conditions on the pilot (or initial) sample size $m( \geq 2)$.
\end{frame}
%%%%%%%%%%%%%%%%%%%%%%%%%%
\begin{frame}{Asymptotic first-order properties}
\vspace{7mm}
For the general MRFWCI estimation methodology carried out under the loss function (5), for every fixed $\alpha, \rho \text{ and }
 \theta$, with $a_i$ defined via assumption $A_i$, we have the following asymptotic first-order conclusions as $d \xrightarrow{} 0$ :

\begin{itemize}
\item[(i)] $n^{*-1}N \xrightarrow[]{P_{\theta}}1$;
\item[(ii)] $E_{\theta}[(n^{*-1}N)^{\kappa}] \xrightarrow[]{} 1$ [\textit{asymptotic first-order efficiency property}];
\item[(iii)] $P_{\theta} \{ \mu \in J_N\} \xrightarrow[]{} (1-\alpha)$ [\textit{asymptotic consistency property}];
\item[(iv)] $\text{RiskEff}_d =\frac{R_N(\mu,J_N)}{R_{n^*}(\mu,J_{n^*})} \xrightarrow[]{} 1$ [\textit{asymptotic risk efficiency property}].
\end{itemize}
 
\end{frame}
%%%%%%%%%%%%%%%%%%%%%%%%%%
\begin{frame}{Asymptotic second-order properties}
\vspace{5mm}
For the general MRFWCI estimation methodology carried out under (5), for every 
fixed $\alpha, \rho \text{ and } \theta$, with $a_i$ defined 
via assumption $A_i$, we have the following asymptotic second-order conclusions 
when $a_3>\frac{5}{2}$ and $a_4=2$ as $d \xrightarrow{} 0$ :
\begin{enumerate}
\item[(i)] $E_{\theta}[(n^{*-1}N)^{1/2}]=1+a_5n^{*-1}+o(n^{*-1})$ where $a_5=\frac{1}{2}(a_1-\frac{1}{2}a_2)$;
\item[(ii)] $P_{\theta} \{ \mu \in J_N\} = (1-\alpha)+2a_6n^{*-1}+o(n^{*-1)}$\\
where $a_6=\frac{1}{2}\{ a_1-(a_2/4)(1-z^2_{\alpha/2}) \}z_{\alpha/2}\phi(z_{\alpha/2})$;
\item[(iii)] $\text{RiskEff}_d =\omega=\frac{R_N(\mu,J_N)}{R_{n^*}(\mu,J_{n^*})}= 1+ a_7 n^{*-1}+o(n^{*-1)}$ \\ 
where $a_7 = 2\{\alpha + 2z_{\alpha/2}\phi(z_{\alpha/2})\}^{-1} \{z_{\alpha/2}\phi(z_{\alpha/2})a_5-a_6\}$.
\end{enumerate}
 
\end{frame}
%%%%%%%%%%%%%%%%%%%%%%%%%%

%%%%%%%%%%%%%%%%%%%%%%%%%%
\section{Sampling Strategies}
%\section{Sampling Strategies}
\begin{frame}{Illustrations of Multistage Sampling Strategies}
    \begin{itemize}
    \item \textbf{Purely Sequential Sampling strategy}
\\ We begin with pilot observations $X_1,\dots ,X_m; m\ge 2$. Then, we record one additional observation 
at-a-time according to the following stopping time:
\begin{equation*}
    N \equiv N_d=inf\{ n \ge m; n \ge z^2_{\alpha/2}S_n^2/d^2 \}  \textrm{ and }  J_N=[\bar{X}_N \pm d].
\end{equation*}

\item \textbf{Accelerated Sequential Sampling strategy}\\
We fix an integer $k(\geq 2)$ and begin with pilot observations $X_1,\dots ,X_m; m\ge 2$. Then, we record 
one additional observation at-a-time according to the following stopping time:        
        $$ T \equiv T_d=inf\{ n \ge m; n \ge k^{-1}z^2_{\alpha/2}S_n^2/d^2 \} $$ 
        $$ N \equiv N_d = kT_d  \text{ and } J_N=[\bar{X}_N \pm d] $$
    \end{itemize}
\end{frame}
%%%%%%%%%%%%%%%%%%%%%%%%%%
\begin{frame}
    \begin{itemize}
 
\item \textbf{Three Stage Sampling strategy}\\
We fix a suitable number $\kappa(\in (0,1))$ and begin with pilot observations $X_1,\dots ,X_m; m\ge 2$ and define:
        $$ T =max\{  m,  <\kappa z^2_{\alpha/2}S_m^2/d^2 >+1\} $$ 
        $$ N =max \{ T,  <z^2_{\alpha/2}S_T^2/d^2 >+1\}  \text{ and } J_N=[\bar{X}_N \pm d] $$
 

\item \textbf{Two-stage Sampling strategy}
    \\ We assume that there exists known $\sigma_L(>0)$ such that $\sigma>\sigma_L$.The pilot sample size is defined as:
    $$m\equiv m_d=max \{ m_0(\ge 2),<z^2_{\alpha/2}\sigma_L^2/d^2>+1\}.$$
    We begin with pilot observations $X_1,\dots ,X_m, m \ge 2$ and find the final sample size as:
    \begin{equation*}
    N \equiv N_d=max\{ m,< t^2_{m-1,\alpha/2}S_m^2/d^2>+1\} \text{ and } J_N=[\bar{X}_N \pm d].
\end{equation*}

    \end{itemize}
\end{frame}
%%%%%%%%%%%%%%%%%%%%%%%%%%

\begin{frame}{}
    \vspace{20mm}
    The table below summarizes some properties of the sampling strategies.
\begin{center}
\begin{tabular}{c c c} 
 \hline
 Sampling Strategy & $a_1$ & $a_2$ \\
 \hline
 Purely Sequential & $-1.1828$ & $2$  \\ 
Accelerated Sequential & $-1.1828k$ & $2k$  \\ 
Three-stage & $0.5-2\kappa^{-1}$ & $2\kappa^{-1}$  \\ 
 Two-stage & $\frac{1}{2}[(z^2_{\alpha/2}+1)\frac{\sigma^2}{\sigma^2_L}+1]$ & $2\frac{\sigma^2}{\sigma^2_L}$ \\
 \hline
\end{tabular}
\end{center}
\end{frame}
%%%%%%%%%%%%%%%%%%%%%%%%%%
\section{Data Analysis and Simulations}
%\section{Data Analysis and Simulations}
\begin{frame}{\textit{Airquality} Data Analysis}

Now we will illustrate applications of the MRFWCI problem for the average in the context of \textit{airquality dataset} 
describing daily air quality measurements in New York. The variable of interest $(X)$ was \textit{wind speed},which
was normally distributed. The mean and standard deviation from the full dataset were: $\mu = 9.957$ and $\sigma = 3.523$. 
The table below summarises the performance of the sampling strategies (fixing $\alpha=0.05$ and $\rho=1.0$). The notations are defined below:
\begin{itemize}
\item $n$: terminal sample size;
\item $\hat{\mu}_n=n^{-1}\sum_{j=1}^n{x_j}$; terminal sample mean;
\item $J_n=[\hat{\mu}_n \pm d]$; terminal $95\%$ MRFWCI for $\mu$.
\end{itemize}
    
\end{frame}

\begin{frame}{}
    \vspace{8mm}
    \begin{center}
\begin{tabular}{c c c c c} 
 \hline
 Sampling Strategy & $d$ & $n$ & $\hat{\mu}_n$ & $J_n$  \\  
 \hline
{Purely Sequential}  & $0.7$ & $102$ & $10.034$ & $[9.334,10.734]$  \\ 
 $m=15$ & $0.8$ & $83$ & $10.427$ & $[9.626,11.226]$ \\
  & $0.9$ & $65$ & $10.795$ & $[9.895,11.695]$\\
 \hline
 {Accelerated Sequential}  & $0.7$ & $88$ & $10.388$ & $[9.687,11.087]$  \\ 
 $m=15, k=2$ & $0.8$ & $74$ & $10.534$ & $[9.733,11.337]$ \\
  & $0.9$ & $56$ & $10.993$ & $[10.093,11.898]$\\
  \hline
{Three-stage }  & $0.7$ & $96$ &$10.214$ & $[9.514,10.914]$  \\ 
 $m=30=5,\kappa=0.5$ & $0.8$ & $72$ & $10.421$ & $[9.612,11.221]$ \\
  & $0.9$ & $59$ & $10.939$ & $[10.039,11.839]$\\
 \hline
\end{tabular}
\end{center}
\end{frame}
%%%%%%%%%%%%%%%%%%%%%%%%%%
\begin{frame}{}
    \vspace{8mm}
    \begin{center}
\begin{tabular}{c c c c c} 
 \hline
 Sampling Strategy & $d$ & $n$ & $\hat{\mu}_n$ & $J_n$  \\  
  \hline
{Two-stage }  & $0.7$ & $105$ &$10.072$ & $[9.375,10.775]$  \\ 
 $\sigma_L=2,m_0=5$ & $0.8$ & $79$ & $10.534$ & $[9.734,11.334]$ \\
  & $0.9$ & $69$ & $10.559$ & $[9.965,11.459]$\\
 \hline
{Two-stage }  & $0.7$ & $115$ & $10.146$ & $[9.446,10.846]$  \\ 
 $\sigma_L=3,m_0=5$ & $0.8$ & $90$ & $10.321$ & $[9.521,11.121]$ \\
  & $0.9$ & $56$ & $10.993$ & $[10.093,11.893]$\\
 \hline
\end{tabular}
\end{center}
\end{frame}

\begin{frame}{Simulation Study}

Simulated performances of  MRFWCI strategies under $10,000$ replications for a 
$N(25,16)$ population with $\alpha=0.05$ and $\rho=1.0$ in (5).
 \begin{center}
\begin{tabular}{c c c c c} 
 \hline
 $n^*$ & $d$ & $\bar{n}$ & $\bar{p}$ & $\bar{\omega}$\\
 \hline
 \multicolumn{5}{c}{Purely Sequential $m=15$}\\
  \hline
 $200$ & $0.5543$ & $198.744$ & $0.9484$  & $1.004$\\ 
 $1000$ & $0.2479$ & $998.048$ & $0.9462$ & $1.001$ \\ 
  $5000$ & $0.1108$ & $4997.536$ & $0.9477$ & $1.001$ \\ 
  \hline
 \multicolumn{5}{c}{Accelerated Sequential $m=15,k=2$}\\
 \hline
 $200$ & $0.5543$ & $196.899$ & $0.9468$ & $1.011$ \\ 
 $1000$ & $0.2479$ & $996.754$ & $0.9505$  & $1.001$\\ 
  $5000$ & $0.1108$ & $4997.327$ & $0.9477$  & $1.001$\\ 
 \hline
\end{tabular}
\end{center}
    
\end{frame}
%%%%%%%%%%%%%%%%%%%%%%%%%%
\begin{frame}{}
    \vspace{10mm}
 \begin{center}
\begin{tabular}{c c c c c} 
 \hline
 $n^*$ & $d$ & $\bar{n}$ & $\bar{p}$ & $\bar{\omega}$\\
 \hline
 \multicolumn{5}{c}{Three-stage $m=15,\kappa=0.5$}\\
  \hline
 $200$ & $0.5543$ & $196.348$ & $0.9457$  & $1.011$\\ 
 $1000$ & $0.2479$ & $995.511$ & $0.9478$ & $1.002$ \\ 
  $5000$ & $0.1108$ & $4995.397$ & $0.9496$ & $1.001$ \\ 
  \hline
\multicolumn{5}{c}{Two-Stage $\sigma_L=3, m_0=10$}\\
 \hline
 $200$ & $0.5543$ & $200.851$ & $0.9459$ & $1.007$ \\ 
 $1000$ & $0.2479$ & $1001.101$ & $0.9459$  & $1.003$\\ 
  $5000$ & $0.1108$ & $4997.536$ & $0.9477$  & $1.001$\\ 
 \hline
\end{tabular}
\end{center}
    
\end{frame}
%%%%%%%%%%%%%%%%%%%%%%%%%%


%%%%%%%%%%%%%%%%%%%%%%%%%%
\section{Conclusion}
\begin{frame}{Conclusion}
\vspace{10mm}
When handling FWCI problems,the lack of a loss function creates a false impression that perhaps 
(i) observations may not cost at all and/or (ii) one's available budget may be unlimited. 
In the MRFWCI formulation, an FWCI problem has been casted in the light of an MRPE by 
balancing the estimation error with the cost of sampling.
\end{frame}
%%%%%%%%%%%%%%%%%%%%%%%%%%
%\section{Reference}

%\nocite{*}
\begin{frame}{Reference}
%\bibliographystyle{unsrt}
%\bibliography{bibfile.bib}

\textbf{Nitis Mukhopadhyay and Swathi Venkatesan}.\\
\vspace{0.2cm}
A New Formulation of Minimum Risk Fixed-Width Confidence Interval
(MRFWCI) Estimation Problems for a Normal Mean with Illustrations
and Simulations: Applications to Airquality Data.\\
\textit{Sequential Analysis, 41, 2022}.

\end{frame}
%%%%%%%%%%%%%%%%%%%%%%%%%%


\end{document}